\section{Harmonic convergence and the stopping-location comparison}\label{app:lemma36-C13}

This appendix gives the full replacement argument for
Lemma~\ref{lemma:convergence-of-harmonic-functions}, the convergence of
discrete harmonic functions and the stopping-location comparison.

\subsection{Definitions and the replacement argument}

Throughout, the reference is the main text.
Fix $\gamma\in(0,2)$, $\rho<1$, and a
deterministic disk $D=B_r(z)\subset B_\rho(0)$.
Take $h$ to be the circle-average embedding of a
$\gamma$-quantum cone, and decorate it by an independent
whole-plane space-filling $\mathrm{SLE}_{16/\gamma^2}$,
subsequently parameterized by $\mu_h$-area and denoted by
$\eta$. Use the associated mated-CRT cell-adjacency graph,
with edge multiplicities counted in simple random walk.
Let $\mathcal F=\sigma(h,\eta)$ and let
$H^\eps_a=\eta([a-\eps,a])$ be the cell of $a\in\eps\mathbb Z$.
Adjacent cells intersect. The cells are connected and locally finite.
Set
\[
 I^\eps_D=\VGeps(D)\setminus\VGeps(\partial D),\qquad
 \sigma^\eps_D=\inf\{j\ge0:X^\eps_j\notin I^\eps_D\},
\]
where $X^\eps$ is one \emph{unrestricted} simple random walk started
at $a_y^\eps=\min\{a\in\eps\mathbb Z:y\in H^\eps_a\}$,
and $\VGeps(A)=\{a:H^\eps_a\cap A\ne\varnothing\}$.
Use this same walk for every stopping operation below.
Write $Y^\eps_j=\eta(X^\eps_j)$ and let $W^\eps$ be its linearly
interpolated embedding, with any positive deterministic time
rescaling, for instance $\meps$ from the main text.
Define
\[
 T_D(W)=\inf\{t\ge0:W(t)\notin D\},\qquad
 Z^\eps=W^\eps(T_D(W^\eps)).
\]
For a path that starts outside $D$ we project its initial point onto
$\overline D$; this harmless convention is irrelevant eventually for
each fixed $y\in D$.

All discrete expressions below are written in terms of the unrestricted walk $X^{y,\eps}$,
and paths that start outside $D$ are projected onto $\overline D$ as above.

For an explicit discontinuity example, take $D=B_1(0)$ and a
real-valued continuous path which runs from $0$ to $1$ in time
$1$, returns to $0$ in the next unit of time, and then runs
at unit speed toward $-\infty$. Its first exit is at $1$.
Multiplying this path by $1-1/n$, $n\ge2$, gives paths converging
locally uniformly to it whose first exit is at $-1$.
For $\psi(w)=\operatorname{Re}w$ the exit functional therefore
does not converge. This is why its limiting-law continuity
set must be checked.

The minimal boundary-value convention is
\[
 \bar\psi=\psi\circ\pi_D,\qquad
 \pi_D:\C\longrightarrow\overline D
\]
the nearest-point projection onto the closed disk.
It is $1$-Lipschitz, so a $\chi$-H\"older constant of $\psi$ is also
one for $\bar\psi$. The minimal path convention for the proof is to
use the unrestricted interpolation $W^\eps$. A precise convention
which also works for the outer-stopped interpolation is
given in Section~\ref{sec:l36-outer}.

\begin{theorem}[Lemma~\ref{lemma:convergence-of-harmonic-functions}, corrected form]\label{thm:l36-A}
For every deterministic $\psi\in C(\overline D)$ and $y\in D$,
\begin{equation}\label{eq:l36-main}
 \mathbb{E}[\bar\psi(Y^\eps_{\sigma^\eps_D})\mid\mathcal F]
 \ \longrightarrow\
 \mathbb{E}[\psi(X^y_{\tau_D})\mid\mathcal F]
 =:u(y)
 \quad\text{in probability}.
\end{equation}
Here $X^y$ is the rescaled LBM in the main text. The function $u$ is
the deterministic classical harmonic extension of
$\psi|_{\partial D}$. In fact the cellwise discrete harmonic
extensions converge uniformly on $\overline D$, in probability.
For H\"older $\psi$ the same-walk comparison in
Theorem~\ref{thm:l36-rate} below gives a polynomial rate.
\end{theorem}

\subsection{Precisely stated inputs}
All page numbers refer to the linked arXiv versions, not journal
pagination.
\begin{enumerate}[label=(\alph*)]
\item Berestycki--Gwynne \cite[Theorem 1.2, p.~10]{berestycki2020random}:
the conditional law of the unrestricted
interpolated embedded walk, with the main-text median time
normalization, converges in probability in the local uniform path
topology to rescaled LBM. The theorem is uniform over starting
points in a fixed compact set.
\item The same paper \cite[Lemma 2.4, p.~17]{berestycki2020random}: for each fixed
$\rho_1<1$ and $0<a<2/(2+\gamma)^2$, with probability
$1-O(\eps^\kappa)$, for some $\kappa>0$, every cell meeting
$B_{\rho_1}$ has diameter at most $\eps^a$.
\item Gwynne--Miller--Sheffield \cite[Lemma 3.11, (3.35), p.~31]{gms-harmonic}:
there are
$\alpha_{\rm ann}=\alpha_{\rm ann}(\gamma)>0$,
$b_{\rm ann}=b_{\rm ann}(\gamma)>0$, and
$p=p(\rho_1,\gamma)>0$ such that, with probability
$1-O(\eps^{\alpha_{\rm ann}})$, simultaneously for every
$c\in B_{\rho_1}$ and
$R\ge\eps^{b_{\rm ann}}$ with $B_{2R}(c)\subset B_{\rho_1}$,
one has
\[
 \inf_{v\in\VGeps(B_R(c))}
 \mathbb{P}_v[\mathcal D_R(c)\mid\mathcal F]\ge p.
\]
Here $\mathcal D_R(c)$ is the event that $X^\eps$ disconnects
$\VGeps(\partial B_{R/2}(c))$ from $\VGeps(\partial B_{2R}(c))$
before hitting the latter set.
Disconnection here means separation in the planar graph by
the vertices visited by the walk. We decrease $p$ to be at
most $1/2$ when convenient.
\item The same paper \cite[Lemma 3.10, (3.34), p.~31]{gms-harmonic}:
for random walk on a graph which almost surely hits a set $A$,
the total variation distance between its first-hit distributions
on $A$, started from $v$ and $w$, is at most
\[
 1-\mathbb{P}_v[X\text{ disconnects }w\text{ from }A
                      \text{ before hitting }A].
\]
We apply this on a finite graph, with extra edges between
stopping vertices if needed to make it connected. The
construction below checks the connectedness and almost-sure
hitting hypotheses explicitly.
\end{enumerate}
The argument below uses these annulus and coupling inputs directly;
it does not require an explicit boundary H\"older exponent.
In particular it avoids a mesh-scale absorption issue in
the proof of Theorem 3.9 of Gwynne--Miller--Sheffield \cite[p.~32, (3.37)]{gms-harmonic}:
uniformly in a boundary distance $u$,
a factor $((\eps^{b_{\rm ann}}\vee s)/u)^c$ cannot in general
be bounded by $C((\eps\vee s)/u)^{c'}$ merely by decreasing
$c'$. For example, if $0<b_{\rm ann}<1$, $u=\eps^{b_{\rm ann}}$
and $s\le\eps$, the first factor is $1$ and the second tends
to $0$. We therefore retain the mesoscopic cutoff and do not
extract numerical comparison exponents from (3.33).

\subsection{The quantitative boundary-excursion estimate}
Fix $\rho<\rho_1<1$ and any
$0<a<2/(2+\gamma)^2$, and set
\begin{equation}\label{eq:l36-exponents}
 \begin{gathered}
 b=\min\{a,b_{\rm ann}\}>0,\qquad d=\eps^a,\qquad
 \ell=100d+\eps^{b_{\rm ann}},\\
 \eps^b\le\ell\le101\eps^b,\qquad
 \nu=\frac{-\log(1-p)}{\log2}>0.
 \end{gathered}
\end{equation}
Choosing $a$ as a function of $\gamma$ makes $b=b(\gamma)$,
whereas $\nu$ can depend on $\rho_1$.
Let $E_\eps\in\mathcal F$ be the intersection of the cell-size event
in (b) with the simultaneous annulus event in (c). Then
\[
 \mathbb{P}(E_\eps^c)\le C\eps^\kappa
\]
for some $\kappa>0$. Notice that $\ell\ge\eps^{b_{\rm ann}}$
and $d<\ell/100$. All annuli used below have radii at least
$\ell$; no estimate is applied below its permitted cutoff.
Only this fixed separation between the cell size and
the first annulus radius is needed.

We first make explicit the planar observation needed with (c).
Let $O$ be a disk, $x\in O$, $B_{3R}(x)\subset B_{\rho_1}$,
and suppose
\[
 R/2>\dist(x,\partial O)+20d.
\]
A walk all of whose visited cells are contained in $O$ cannot
disconnect $\VGeps(\partial B_{R/2}(x))$ from
$\VGeps(\partial B_{2R}(x))$.
Indeed take the outward radial segment about the center of
$O$ from $u\in\partial B_{R/2}(x)$ to
$w\in\partial B_{2R}(x)$. Its distance from $O$ is greater
than $20d$. Choose cells containing $u,w$; their vertices
belong to the respective circle vertex sets.
Use the usual planar cell-adjacency drawing, with vertices
in their cells and edges in arbitrarily small neighborhoods
of the unions of incident cells. Choose perturbations of
size less than $d/10$ to accommodate pinch points.
Join the two chosen embedded vertices to $u,w$, respectively,
and insert the radial segment between these short joins.
The resulting curve stays at distance greater than $18d$
from $O$.

Here is why following faces along this curve is legitimate
even when cells have nontrivial topology. In the buffered
region, every point is within $2d$ of an embedded vertex,
since some cell contains the point. Consequently a face
meeting the curve has boundary within $2d$ of it; otherwise
that nearby vertex would lie in the interior of the face.
An edge there belongs to neighborhoods of its incident cells.
These cells, and the third incident cell of the triangular
face, all meet the ambient cell-size region. Each has
diameter at most $d$. They intersect pairwise, so their
union has diameter at most $3d$; the face boundary, and
hence the bounded face itself, have diameter less than $4d$.
The latter assertion follows because a bounded face is
contained in the convex hull of its boundary.

Local finiteness now gives a finite chain of closed faces
along the curve, consecutive faces sharing an edge or a
vertex. Their boundary vertices give a graph path between
the two chosen vertices. Every cell on this path is within
$6d$ of the curve and is therefore disjoint from $O$.
Thus the path avoids every visited vertex, contradicting
graph disconnection. This also checks the endpoint
membership in the two prescribed circle vertex sets.
Equivalently this is the planar separation
observation in the proof of GMS Theorem 3.9 \cite[p.~32, immediately
before (3.38)]{gms-harmonic}; the buffer here makes its application insensitive
to cell representatives. This auxiliary planar drawing is
used only for graph separation; $W^\eps$ remains the
straight-line interpolation specified at the start.

\begin{lemma}[Excursion after a boundary-cell hit]\label{lem:l36-excursion}
On $E_\eps$, for all sufficiently small $\eps$, uniformly
over walks started in $I_D^\eps$,
\begin{equation}\label{eq:l36-spatial-tail}
 \mathbb{P}\bigl[|Z^\eps-Y^\eps_{\sigma_D^\eps}|>C s\mid\mathcal F\bigr]
 \le \min\{1,C(\ell/s)^\nu\},\qquad 0<s\le s_0,
\end{equation}
where $s_0>0$ is deterministic. The same estimate bounds the
maximal displacement from $Y^\eps_{\sigma_D^\eps}$ after that
visit and before Euclidean exit, if exit has not already
occurred on the entering edge.
\end{lemma}
\begin{proof}
An interior cell is contained in $D$, by connectedness.
The first noninterior cell reached from $I_D^\eps$ meets
$\partial D$: it intersects its preceding interior cell and
cannot be wholly outside $D$. Its representative is therefore
within $d$ of $\partial D$.
Before the edge entering it, all interpolated segments lie in
the convex disk $D$. If Euclidean exit occurs on the entering
edge, its distance from the endpoint is at most $2d$, since
adjacent cell representatives have distance at most $2d$.

Otherwise let $v=X^\eps_{\sigma_D^\eps}$,
$x=\eta(v)\in D$, and put $O=B_{r+3d}(z)$.
The cell of $v$ is contained in $O$, and
$\dist(x,\partial O)\le4d$.
A first visit to $\VGeps(\partial O)$ has representative at
distance at least $r+2d$ from $z$.
Consequently the interpolation has exited $D$ by that visit.
It is enough to control the walk's excursion before hitting
$\VGeps(\partial O)$.

Let $R_j=2^j\ell$, and retain indices with $2R_j\le s_0$,
where $s_0<(\rho_1-\rho)/2$ is small enough that even
$B_{3R_j}(x)\subset B_{\rho_1}$. At the initial stage
the walk is at $v$.
At each subsequent stage restart it at its first hit of
$\VGeps(\partial B_{2R_{j-1}}(x))$, unless it has already hit
$\VGeps(\partial O)$. The restart cell meets the open disk
$B_{R_j}(x)$ even though $2R_{j-1}=R_j$:
it intersects the preceding cell, which is contained
in $B_{R_j}(x)$ since this is the first boundary-cell hit.
Its representative need not lie in that open disk, but
input (c) requires exactly membership in $\VGeps(B_{R_j}(x))$.
Also the restart cell cannot already meet
$\partial B_{2R_j}(x)$, since its diameter is at most
$d<R_j/100$. Thus consecutive dyadic radii give legitimate
successive stopping times.
For small $\eps$ the planar observation applies at every
stage, because $R_j/2\ge\ell/2>50d$ whereas
$\dist(x,\partial O)+20d\le24d$.
Thus the disconnection event in input (c) forces a visit
to $\VGeps(\partial O)$ before reaching
$\VGeps(\partial B_{2R_j}(x))$.
The conditional chance of surviving to that outer circle
without visiting $\VGeps(\partial O)$ is at most $1-p$.

The strong Markov property, applied successively to these
stopping times, now gives a survival probability at most
$(1-p)^{j+1}$ through stage $j$. The annuli overlap; no
independence is being used. The uniform conditional bound
$1-p$ at each restart is the only input to the product.
Before the first boundary-cell hit for $B_{2R_j}(x)$,
all representatives lie in $B_{2R_j}(x)$; the representative
of the hit cell lies in $B_{2R_j+d}(x)$. Convexity gives the
same bound for interpolated segments.
For a displacement scale $s\ge4\ell$, choose the largest $j$
with $2R_j\le s/2$. Since $R_j=2^j\ell$,
\[
 (1-p)^{j+1}\le C(\ell/s)^\nu.
\]
A displacement greater than a fixed multiple of $s$ before
Euclidean exit requires survival through this stage.
The cases $s<4\ell$ follow by enlarging the constant and using
the bound $1$. Conditional on the environment the annulus
event is uniform over all centers, radii, and starting
vertices, so it remains applicable to the random $x$.
Finally apply the strong Markov property at $\sigma_D^\eps$.
\end{proof}

\begin{theorem}[Quantitative comparison]\label{thm:l36-rate}
Suppose $\psi$ is $\chi$-H\"older on $\overline D$, with
$0<\chi\le1$ and H\"older constant $H$.
Fix $y\in D$ and start the walk at $a_y^\eps$.
Define $\Phi^\psi(W)=\psi(W(T_D(W)))$ on paths which start
inside $D$ and exit in finite time, and set it equal to $0$
on all remaining paths. For every
\[
                         0<q<b\chi
\]
one may take $C_\psi=1+H$ so that, for small $\eps$,
\begin{equation}\label{eq:l36-rate}
 \mathbb{P}\!\left[
 |\Phi^\psi(W^\eps)-\bar\psi(Y^\eps_{\sigma_D^\eps})|
       >C_\psi\eps^q\right]
 \le C\bigl(\eps^\kappa+\eps^{\nu(b-q/\chi)}\bigr).
\end{equation}
On $E_\eps$ the corresponding conditional probability is at
most $C\eps^{\nu(b-q/\chi)}$, uniformly over interior starting
vertices. Moreover, on $E_\eps$,
\begin{equation}\label{eq:l36-mean-rate}
 \mathbb{E}[|\Phi^\psi(W^\eps)-\bar\psi(Y^\eps_{\sigma_D^\eps})|
       \mid\mathcal F]\le
 \begin{cases}
 C H\eps^{b\chi},&0<\chi<\nu,\\
 C H\eps^{b\nu}(1+|\log\eps|),&\chi=\nu,\\
 C H\eps^{b\nu},&\nu<\chi\le1 .
 \end{cases}
\end{equation}
In particular, choosing $a=(2+\gamma)^{-2}$ permits
\[
 q=\frac{\chi}{2}\min\{(2+\gamma)^{-2},b_{\rm ann}(\gamma)\},
\]
so that $q$ depends only on $\chi$ and $\gamma$.
More generally, by choosing $a$ close enough to its upper
endpoint, every
\begin{equation}\label{eq:l36-full-q-range}
 0<q<\chi\min\{2/(2+\gamma)^2,b_{\rm ann}(\gamma)\}
\end{equation}
is allowed.
\end{theorem}
\begin{proof}
For fixed $y\in D$, its starting cell is interior on $E_\eps$
for all small $\eps$.
The unrestricted walk exits the finite cell domain associated
with $B_{r+3d}(z)$ almost surely, so its interpolation exits $D$.
Projection is nonexpansive and fixes its exit point
$Z^\eps\in\partial D$. Hence
\[
 |\psi(Z^\eps)-\bar\psi(Y^\eps_{\sigma_D^\eps})|
 \le H|Z^\eps-Y^\eps_{\sigma_D^\eps}|^\chi.
\]
Use Lemma~\ref{lem:l36-excursion} with $Cs=\eps^{q/\chi}$,
which is larger than a fixed multiple of $\ell$ for small
$\eps$, and add $\mathbb{P}(E_\eps^c)$.
For the mean, write $R=|Z^\eps-Y^\eps_{\sigma_D^\eps}|
\le2r+d$. By monotonicity the tail bound extends, with a
different constant, to $s\le2r+d$.
Split
\[
 \mathbb{E}[R^\chi\mid\mathcal F]
 =\chi\int_0^{2r+d}s^{\chi-1}\mathbb{P}[R>s\mid\mathcal F]\,ds
\]
at a constant multiple of $\ell$.
The remaining integral is bounded by
$C\ell^\nu\int_\ell^{2r+d}s^{\chi-\nu-1}\,ds$.
Its three elementary cases give \eqref{eq:l36-mean-rate}.
\end{proof}
The ambient radius enters $\nu$, the probability exponent
$\kappa$, and multiplicative constants, but not the choice
$q=b(\gamma)\chi/2$ or $C_\psi=1+H$.
This is a rate comparing two stopping rules for the same walk.
It supplies neither a convergence rate to Brownian harmonic
measure nor an optimal numerical exponent.

\subsection{Continuity at the limiting paths and proof of the lemma}
\begin{lemma}[Exit-map continuity]\label{lem:l36-exit-continuity}
Let $w$ start in $D$, let $\tau=T_D(w)<\infty$, and suppose that
for each $\delta>0$ some $t\in(\tau,\tau+\delta)$ has
$w(t)\notin\overline D$. At such a path the exit time and exit
location maps are continuous for local uniform convergence.
\end{lemma}
\begin{proof}
For each $s<\tau$, the compact set $w([0,s])$ lies at positive
distance from $D^c$, so all sufficiently close paths have exit
time greater than $s$. Choose $t\in(\tau,\tau+\delta)$ at positive
distance from $\overline D$; sufficiently close paths exit by $t$.
Thus their exit times converge to $\tau$. Uniform continuity of
$w$ on a compact time interval containing those times then proves
convergence of the exit locations.
\end{proof}

For Brownian motion the exit time from a disk is finite almost
surely. At its exit point $p$, projection of the subsequent Brownian
increments onto the outward unit normal is a one-dimensional
Brownian motion. The probability that this projection stays
nonpositive on any interval $(0,\delta)$ is zero (for example,
the reflection principle gives the distribution of its maximum).
The strong Markov property and a countable intersection over
rational $\delta>0$ imply the crossing hypothesis of
Lemma~\ref{lem:l36-exit-continuity}. LBM is the same Brownian trace
under a continuous strictly increasing clock, finite up to exit
from bounded disks, so it has the same property.
These assertions hold for almost every environment under its
conditional LBM law, by disintegration. Time rescaling by $m_0$
does not change the exit location.

For the path-law argument choose the median clock $\meps$
specified in input (a). All stopping locations and values
under consideration are unchanged by positive deterministic
time rescaling, so this choice does not restrict the result.
Now start with an arbitrary deterministic sequence
$\eps_j\to0$. By input (a), extract a deterministic subsequence
along which the conditional path laws converge almost surely.
For each environment in this event the bounded exit functional
is continuous almost everywhere under the limiting law.
The bounded almost-everywhere-continuous version of the
portmanteau theorem therefore gives
\[
 \mathbb{E}[\Phi^\psi(W^\eps)\mid\mathcal F]
 \longrightarrow\mathbb{E}[\psi(X^y_{\tau_D})\mid\mathcal F]
 \quad\text{almost surely along that subsequence}.
\]
Here one may see the portmanteau implication directly by replacing a
bounded function by its upper and lower $\delta$-oscillation
envelopes: the envelopes converge to the function at continuity
points, and the upper and lower weak-convergence inequalities,
followed by bounded convergence, squeeze their integrals.
No assertion that the exit functional is continuous everywhere
has been used.

For H\"older $\psi$, \eqref{eq:l36-mean-rate} and
$\mathbb{P}(E_\eps^c)\to0$ show that replacing
$\Phi^\psi(W^\eps)$ by $\bar\psi(Y^\eps_{\sigma_D^\eps})$
changes its conditional expectation by a quantity tending to
zero in probability. The arbitrary-subsequence criterion now
proves \eqref{eq:l36-main}. Equivalently, choose a further
deterministic subsequence with summable exceptional probabilities.

The LBM exit distribution is the Brownian exit distribution.
The latter has density
\[
 u(y)=\int_0^{2\pi}
 \psi(z+re^{it})\,
 \frac{r^2-|y-z|^2}{|re^{it}-(y-z)|^2}\,\frac{dt}{2\pi}.
\]
The Poisson kernel is harmonic in $y$, integrates to $1$, and
concentrates at the boundary point as $y$ tends to that point.
This proves harmonicity, continuity up to the boundary, and the
assertion that $u$ is deterministic.

We prove the uniform assertion directly, also keeping the
cutoff in the compactness argument.
Write $f^\eps(v)=\mathbb{E}_v[\bar\psi(Y^\eps_{\sigma_D^\eps})\mid\mathcal F]$,
so $\|f^\eps\|_\infty\le\|\psi\|_\infty$.
Here is the finite graph to which input (d) is applied.
Keep $I_D^\eps$ and the stopping vertices
$A=\VGeps(\partial D)$, together with all original edges incident
to $I_D^\eps$. Every neighbor of an interior vertex belongs
to $I_D^\eps\cup A$, by connectedness of its cell.
Every component containing an interior vertex reaches $A$:
otherwise it would be a finite component of the connected
infinite mated-CRT graph.
Add edges between vertices of $A$ to connect the remaining
components if necessary. This finite graph is connected,
and its walk hits $A$ almost surely from every vertex.
The walk before that hit, its hitting distribution, and
the event of disconnecting a vertex from $A$ are unchanged
by the added edges. Thus all hypotheses of input (d) hold.
Fix a compact $K\Subset D$ and let
$\delta=\dist(K,\partial D)>0$.
For representatives $x=\eta(v)$ and $x'=\eta(w)$ near $K$,
start the annulus iteration at
\[
 R_0=8(\ell\vee|x-x'|)
\]
and use radii $2^jR_0$ until $2\cdot2^jR_0\le\delta/4$.
If $R_0$ is already larger than this bound the following
estimate is trivial after enlarging its constant.
Otherwise the cell of $w$ lies inside the innermost circle,
and every cell of $\partial D$ lies beyond the outermost
circle, for small $\eps$. A disconnection at any stage
therefore separates $w$ from the stopping set
$\VGeps(\partial D)$. The walk cannot hit that stopping set
before hitting the stage's outer circle.
The Markov iteration in Lemma~\ref{lem:l36-excursion}, now counting
failure of all disconnection attempts, bounds its probability
by $C(R_0/\delta)^\nu$.
Input (d) and the expectation formula consequently give
\begin{equation}\label{eq:l36-interior-modulus}
 |f^\eps(v)-f^\eps(w)|
 \le C\|\psi\|_\infty
       \left(\frac{\ell\vee|x-x'|}{\delta}\right)^\nu.
\end{equation}
All constants here are uniform in the vertices of the
chosen compact set.

For boundary control, take an interior representative $x$
at distance $u$ from $\partial D$ and start the same iteration
at $R_0=8(\ell+u)$, now stopping on $\VGeps(\partial D)$ itself.
For $u$ small this gives
\[
 \mathbb{P}_x[|Y^\eps_{\sigma_D^\eps}-x|>Cs\mid\mathcal F]
 \le \min\{1,C((\ell+u)/s)^\nu\},\qquad 0<s\le s_0.
\]
The proof is identical to Lemma~\ref{lem:l36-excursion}:
the radial complement of $D$ forces a boundary hit on each
successful disconnection, and $R_0/2>u+20d$.
Let $x_\partial$ be the nearest point of $\partial D$ to $x$.
With
$\omega_\psi(t)=\sup\{|\psi(v)-\psi(w)|:
v,w\in\overline D,\ |v-w|\le t\}$,
nonexpansiveness of projection implies
\begin{equation}\label{eq:l36-boundary-modulus}
 |f^\eps(v)-\psi(x_\partial)|
 \le \omega_\psi(Cs+u)
       +C\|\psi\|_\infty((\ell+u)/s)^\nu .
\end{equation}
Choose $s=\sqrt{\ell+u}$ for small $\ell+u$.
The right side tends to zero uniformly as $\ell+u\to0$.
For a starting boundary cell the prescribed value differs
from the value at a boundary point in the cell by at most
$\omega_\psi(d)$, so the same limiting assertion holds.
This includes cells touching the boundary only from outside.

Now prepare, from any deterministic sequence of $\eps$'s,
a deterministic subsequence on which convergence holds at a
countable dense set of points in $D$ and $E_\eps$ holds
eventually almost surely. Diagonal selection with summable
error probabilities does this.
On a compact subset of $D$, \eqref{eq:l36-interior-modulus} and
$|\eta(a_y^\eps)-y|\le d$ give asymptotic equicontinuity of
the cellwise extensions. A finite net from the dense set
therefore upgrades their pointwise convergence to uniform
convergence on that compact subset.
Estimate \eqref{eq:l36-boundary-modulus} controls the remaining
boundary strip, uniformly in $\eps$, and the limiting
Poisson extension is uniformly continuous on $\overline D$.
First shrink the strip and then let $\eps\to0$.
This proves almost-sure uniform convergence along the
prepared subsequence, hence uniform convergence in probability
by the arbitrary-subsequence criterion.

Finally approximate a continuous datum uniformly on
$\overline D$ by deterministic Lipschitz functions. Projection
preserves the uniform approximation, and both discrete and
continuum harmonic extension are contractions in the sup norm,
by their expectation formulas. The uniform conclusion and
\eqref{eq:l36-main} follow for continuous $\psi$ as well.
This proves Theorem~\ref{thm:l36-A}.

\subsection{The actual outer-stopped interpolation}\label{sec:l36-outer}
Here is a precise repair if one wants to retain
$\widehat X^{y,\eps}$, the walk stopped on outer boundary cells. For an eventually
constant path $w$ which starts in $D$, set
\[
 \Phi^\psi_{\rm stop}(w)=
 \begin{cases}
 \psi(w(T_D(w))),&T_D(w)<\infty,\\
 \bar\psi(w(\infty)),&T_D(w)=\infty.
 \end{cases}
\]
Use the projected initial value when the path starts outside $D$.
With this convention \eqref{eq:l36-rate} and \eqref{eq:l36-mean-rate}
hold with $\Phi^\psi_{\rm stop}(\widehat X^{y,\eps})$ in place
of $\Phi^\psi(W^\eps)$. Indeed $\sigma_D^\eps\le
\sigma_{B_\rho}^\eps$. If the outer-stopped path has not exited
$D$, its final point occurs during the excursion before the
unrestricted path's first exit. The maximal-excursion assertion
of Lemma~\ref{lem:l36-excursion} bounds its displacement from
$Y^\eps_{\sigma_D^\eps}$ by the same argument.
Projection can only decrease the relevant distance.

If $\overline D\subset B_\rho$, the outer-stopped path does exit
$D$ on the cell-size event for small $\eps$, and this convention
agrees with the functional $Z^\eps$ defined above. Without such a separation,
one must supply a convention or use the unrestricted walk.
For the proof of harmonic convergence, the unrestricted walk
avoids any need to apply a continuity theorem to this extended
outer-stopped functional.

For any two fixed H\"older extensions of the same $\psi$, their
values at a boundary-cell representative differ by $O(d^\chi)$:
compare each to a point of $\partial D$ in that cell.
Thus the limiting lemma is independent of the chosen extension.
With the conventions above, the polynomial comparison holds for the unrestricted walk.
